\documentclass[10pt,a4paper]{article}

\usepackage[margin=2.0cm]{geometry}
\usepackage[utf8]{inputenc}
\usepackage[T1]{fontenc}
\usepackage{lmodern}
\usepackage{microtype}
\usepackage{booktabs}
\usepackage{tabularx}
\usepackage{ragged2e}
% Ragged-right X column so the line-breaker doesn't have to flush-right
% texttt-heavy descriptions (which overflow when forced to justify).
\newcolumntype{Y}{>{\RaggedRight\arraybackslash}X}
\usepackage{hyperref}
\usepackage{xcolor}
\usepackage{enumitem}
\usepackage{fancyhdr}

% Allow LaTeX a little slack when line-breaking inside tabularx X columns
% with many \texttt{} fragments (unbreakable strings).
\setlength{\emergencystretch}{3em}

\hypersetup{
  colorlinks=true,
  linkcolor=black,
  urlcolor=blue!50!black,
  citecolor=black,
}

\definecolor{cgreen}{HTML}{1A7F37}
\definecolor{cred}{HTML}{C9302C}
\definecolor{camber}{HTML}{B58105}

\setlist[itemize]{leftmargin=*,topsep=2pt,itemsep=1pt}

\pagestyle{fancy}
\fancyhf{}
\fancyhead[L]{\small Gate detection status}
\fancyhead[R]{\small ETHRC Aeronautics --- 2026-05-10}
\fancyfoot[C]{\small \thepage}
\renewcommand{\headrulewidth}{0.3pt}

\title{\vspace{-1.5em}Gate detection --- what works \& what is left to test}
\author{Mahbod Tajdini \\ \small \texttt{mtajdini@ethrobotics.ch} --- ETHRC Aeronautics}
\date{2026-05-10}

\begin{document}
\maketitle
\thispagestyle{fancy}
\vspace{-1em}

\section*{TL;DR}

The detection pipeline (segmentation + multi-task keypoint head + classical
QuAdGate) trains, exports to TensorRT, and produces sub-pixel-accurate corners
on TII piloted footage. The deployment-headline checkpoint
\texttt{best\_deeplabv3plus\_keypoint.pth} achieves \textbf{85.1\,\% detection rate}
at \textbf{1.01\,px median} corner error on flight-01p-ellipse, single forward
pass, no prior, no EKF state, 102\,fps wall-clock on an RTX 5090 (0.47\,ms TRT
FP16). What is \emph{not} yet measured: detection rate on flights other than
the one we keep benching, false-positive rate on no-gate frames, real-recorded
footage (blocked on Drive data).

\section{What works today}

Numbers are on \texttt{flight-01p-ellipse} (TII v3.0.0 piloted, 423 frames,
seed 42), reproducible from
\href{run:../07_keypoint_head/findings.md}{report 07}.

\begin{description}[leftmargin=1.4em,style=nextline,itemsep=2pt,topsep=2pt]
  \item[Segmentation backbone.]
    DeepLabV3+ with ResNet-34 + ImageNet pretrain --- val IoU \textbf{0.872}.
    \hfill {\small\itshape report 05 / 07}
  \item[Keypoint head.]
    Multi-task 4-channel Gaussian-heatmap head, sub-pixel decode via
    soft-argmax. \hfill {\small\itshape report 07}
  \item[Deployment-mode bench (\texttt{model\_direct}).]
    \textbf{85.1\,\% detection} / \textbf{1.01\,px median} / 63.7\,\% within 5\,px
    on flight-01p-ellipse (no GT prior, no EKF state).
    \hfill {\small\itshape report 07}
  \item[Prior-mode framework.]
    Five priors wired in one bench script: \texttt{gt}, \texttt{noisy},
    \texttt{none}, \texttt{model}, \texttt{model\_direct}, \texttt{ekf}.
    \hfill {\small\itshape reports 06 / 07 / 08}
  \item[TRT FP16 export.]
    Multi-task graph exports to ONNX (opset 18) and builds an FP16
    TRT engine in $\sim 6\,$s. \textbf{0.47\,ms / 2128 fps on RTX 5090},
    projects $\sim$50--100\,fps on Jetson Orin NX.
    \hfill {\small\itshape report 03 + \texttt{bench\_trt\_keypoint.py}}
  \item[TRT numerical correctness.]
    TRT FP16 output matches torch FP32 to
    $|trt - torch|_\infty < 5\!\times\!10^{-4}$ on uniform-$[0,1]$ input.
    \hfill {\small\itshape same}
  \item[Multi-gate sibling checkpoint.]
    Multi-gate Gaussian targets + coherent-quad decode at inference;
    88.7\,\% det / 1.26\,px median (tied with report 07 on within-5\,px).
    \hfill {\small\itshape report 10}
\end{description}

\paragraph{Knock-on consequences ruled in.}
TII-distribution detection is solved (within the constraints of the
report-07 metric); architecture exports cleanly to TRT; multi-task did
not regress seg val IoU; the multi-gate-target experiment closed without
moving the headline either way.

\section{What is \emph{not} measured (and matters for detection reliability)}

\begin{tabularx}{\linewidth}{@{}c Y r@{}}
\toprule
\textbf{\#} & \textbf{Open test} & \textbf{Cost} \\
\midrule
1 & Cross-flight bench: pick 4--5 flights from \texttt{*-ellipse}, \texttt{*-lemniscate}, \texttt{*-trackRATM}. Tells us if 85\,\% generalises across camera trajectory and track layout. & {\small \textcolor{cgreen}{30 min}} \\
2 & False-positive rate on no-gate frames: scan the dataset for empty label files (start of each flight, between-gate transits), measure how often the model fires anyway at the current 0.3 threshold. & {\small \textcolor{cgreen}{30 min}} \\
3 & Cross-flight bench of the multi-gate sibling ckpt: does the +3.6\,pp detection-rate gain in report 10 hold out of distribution? & {\small \textcolor{cgreen}{15 min}} \\
4 & Confidence calibration: replace $\text{peak} \!\geq\! 0.3$ with peak~+~spatial concentration (KL divergence of the peak's $7\!\times\!7$ window from a $\sigma\!=\!3$ Gaussian). Catches doorway-style false positives in report 07's sample panel. & {\small \textcolor{camber}{2 h}} \\
5 & Bench on \texttt{flight-01p-fixed} (the original 82\,\% headline): blocked on Drive data + missing \texttt{scripts/prepare\_labels.py}. & {\small \textcolor{cred}{blocked}} \\
6 & Real-footage bench (the README's 39\,\% number): blocked on Drive data. & {\small \textcolor{cred}{blocked}} \\
7 & Heavier motion-blur augmentation + retrain. Current aug caps at 15\,px blur kernel; trackRATM samples wander on fast-turn frames. & {\small \textcolor{camber}{1 day}} \\
8 & Real-data fine-tuning (encoder frozen, decoder + head on team-recorded footage). & {\small \textcolor{cred}{blocked}} \\
\bottomrule
\end{tabularx}

\paragraph{What the open tests do \emph{not} cover.}
``Which gate is the next one to fly through'' is not a perception question
under the current architecture --- it is a planner-side concern that needs
an EKF + track-plan + plane-test (the MonoRace recipe, Bahnam et al. 2026).
That workstream is items 4--6 on the master rank list (Kalman fusion,
ROS2, Orin NX) and intentionally outside the scope of ``are gates
detected reliably.''

\section{Recommended order}

Items 1 and 2 cost about an hour combined and are pure measurement. They
tell us whether items 4 / 7 are actually worth doing, or whether the
current numbers already generalise. Recommended sequence:

\begin{enumerate}
  \item[\textbf{1.}] \textbf{Item 1 (cross-flight bench)} first --- tells us
        whether 85.1\,\% is representative or one-flight luck. ~30 min.
  \item[\textbf{2.}] \textbf{Item 2 (false-positive rate)} --- closes the
        precision picture. If FP rate is already $<\!2\,\%$, item 4 isn't
        worth doing. ~30 min.
  \item[\textbf{3.}] \textbf{Item 4 (confidence calibration)} only if item
        2 shows a real false-positive problem. ~2 h.
  \item[\textbf{4.}] \textbf{Item 3 (multi-gate sibling cross-flight)} ---
        informational, can be skipped without affecting the deployment path.
\end{enumerate}

Items 5--6 unblock once the Drive data lands on qolam; item 7 is the only
remaining unblocked retrain lever and only justifies its cost if item 1
shows trajectory-specific failures.

\paragraph{Bottom line.}
The model detects gates well on the slice we have benched, and we have
enough plumbing (TRT export, prior-mode framework, multiple checkpoints)
to ship something. The remaining ``make sure'' work is mostly
empirical measurement on un-benched slices --- about 1\,h of work, no
retraining, no code changes for items 1--3.

\end{document}
